\chapter{其他线性注意力变体}

\section{概述}

在 Linear Transformer 和 Performer 之后，研究者提出了多种线性注意力的变体和改进方法，旨在提高表达能力、数值稳定性和训练效果。本章介绍几种重要的变体。

\section{MGA (Memory-based Global Attention)}

MGA 结合了线性注意力和记忆机制，通过引入一个小型的"记忆矩阵"来增强线性注意力的表达能力。

\begin{equation}
    \mathbf{o}_i = \text{LinearAttn}(\mathbf{q}_i, \{\mathbf{k}_j, \mathbf{v}_j\}) + \mathbf{M}^\top \mathbf{q}_i
\end{equation}

其中 $\mathbf{M} \in \mathbb{R}^{d \times r}$ 是可学习的记忆矩阵，$r \ll N$ 是记忆槽的数量。这种方法允许模型在保持线性复杂度的同时，学习全局的模式表示。

\section{GLA (Gated Linear Attention)}

GLA 将门控机制引入线性注意力，以增强对信息的控制能力。其核心思想是：

\begin{equation}
    \mathbf{S}_t = \alpha_t \odot \mathbf{S}_{t-1} + \beta_t \odot (\phi(\mathbf{k}_t)\mathbf{v}_t^\top)
\end{equation}

其中 $\alpha_t, \beta_t \in \mathbb{R}^{d \times d}$ 是门控信号，由查询和输入动态生成。$\odot$ 表示逐元素乘法。

门控机制允许模型学习何时更新状态、何时保留历史信息，类似于 LSTM 中的遗忘门和输入门。

\section{RWKV}

RWKV（Receptance Weighted Key Value）是一种将 RNN 效率和 Transformer 质量相结合的架构。它本质上是一种特殊的线性注意力，其递推形式为：

\begin{equation}
    \mathbf{wkv}_t = \text{diag}(\mathbf{u}) \cdot \mathbf{v}_t + \sum_{i=1}^{t-1} \text{diag}(\mathbf{w})^{t-1-i} \cdot \mathbf{v}_i
\end{equation}

其中 $\mathbf{u}$ 和 $\mathbf{w}$ 是可学习的时间衰减参数。RWKV 的关键创新在于：

\begin{itemize}
    \item 使用通道混合（channel mixing）替代前馈网络
    \item 通过时间衰减 $\mathbf{w} \in (0, 1)$ 实现类似 RNN 的遗忘机制
    \item 训练时通过并行扫描（parallel scan）实现并行化
    \item 推理时是纯 RNN 形式，$\mathcal{O}(1)$ 每步
\end{itemize}

\section{Retentive Network (RetNet)}

RetNet 由 Microsoft 于 2023 年提出，核心思想是在训练时使用并行形式，在推理时使用递归形式。

RetNet 使用 multi-scale decay 来替代标准注意力中的 softmax：

\begin{equation}
    \mathbf{A}_{i,j} = \gamma^{i-j} \cdot \exp(\mathbf{q}_i^\top \mathbf{k}_j / \sqrt{d}), \quad j \leq i
\end{equation}

其中 $\gamma \in (0, 1)$ 是衰减因子。这种设计使得：

\begin{itemize}
    \item 训练时：可以使用累积和并行计算所有位置的输出
    \item 推理时：状态递推更新，每步 $\mathcal{O}(d^2)$
\end{itemize}

\section{HGRN (Hierarchically Gated RNN)}

HGRN 结合了线性注意力和层次门控机制。其递推公式为：

\begin{equation}
    \mathbf{S}_t = \mathbf{g}_t \odot \mathbf{S}_{t-1} + \tilde{\mathbf{k}}_t \tilde{\mathbf{v}}_t^\top
\end{equation}

其中 $\mathbf{g}_t = \sigma(\mathbf{W}_g \mathbf{x}_t)$ 是门控信号。HGRN 通过多层堆叠实现了从局部到全局的信息整合。

\section{基于卷积的线性注意力}

一些方法将线性注意力与卷积操作结合。基本思想是：

\begin{equation}
    \mathbf{o}_i = \sum_{j=1}^N \kappa(i-j) \cdot \phi(\mathbf{q}_i)^\top \phi(\mathbf{k}_j) \mathbf{v}_j
\end{equation}

其中 $\kappa(i-j)$ 是基于位置距离的衰减函数。这种方法引入了位置感知，弥补了纯线性注意力对位置信息敏感不足的问题。

\section{变体比较}

\begin{table}[H]
    \centering
    \begin{tabular}{l|cccc}
        \toprule
        方法 & 复杂度 & 递推 & 门控 & 并行训练 \\
        \midrule
        Linear Transformer & $\mathcal{O}(Nd^2)$ & \checkmark & \ding{55} & \checkmark \\
        Performer & $\mathcal{O}(Nmd)$ & \checkmark & \ding{55} & \checkmark \\
        GLA & $\mathcal{O}(Nd^2)$ & \checkmark & \checkmark & \checkmark \\
        RWKV & $\mathcal{O}(Nd)$ & \checkmark & \checkmark & \checkmark \\
        RetNet & $\mathcal{O}(Nd^2)$ & \checkmark & \ding{55} & \checkmark \\
        \bottomrule
    \end{tabular}
    \caption{线性注意力变体特性比较}
\end{table}
